% !TeX root = ../教材题库_全部条目_转写工作稿.tex
% Source: XJ-B1-C03-HW-P085-Y1280
% 第3章 / 3.2 函数的基本性质 / 教材 p.85、86 / PDF 第92、93页
% OCR未校对：公式与文字均可能有误；未完成数学TeX转写；禁止直接发给学生
% 原题及相邻原解截图（核对用）：../题目原图/XJ-B1-C03-HW-P085-Y1280-01.webp, ../题目原图/XJ-B1-C03-HW-P085-Y1280-02.webp
% 本文件只在转写完成后作为题干使用；OCR 条目尚不是可靠题干。
\begin{bankitem}{习题3.2 第5题}{第3章 / 3.2 函数的基本性质 / 教材 p.85、86 / PDF 第92、93页}{XJ-B1-C03-HW-P085-Y1280}{OCR未校对：公式与文字均可能有误；未完成数学TeX转写；禁止直接发给学生}
\begin{ocrdraft}
[教材p.85 / PDF第92页]
\par
5．已知不与g(x）分别是定义在R上的偶函数和奇函数，将下图补充完整．
\par
g0）
\par
丆0）
\par
0
\par
0
\par
（第5题）
\par
[教材p.86 / PDF第93页]
\end{ocrdraft}
\end{bankitem}
